\documentclass{article}
\usepackage{shortex}
\usepackage[margin=1in]{geometry}
\usepackage{listings}
\usepackage{xcolor}
% \usepackage{jupynotex}
\usepackage{csvsimple}
\usepackage{placeins}
% \usepackage{hyperref}


\graphicspath{{./graphics}}

\lstdefinelanguage{Julia}{
    keywords={function, end, if, else, elseif, for, while, return, 
               local, global, const, struct, mutable, abstract, type,
               using, import, export, module, begin, let, do, try, 
               catch, finally, true, false, nothing},
    sensitive=true,
    comment=[l]{\#},
    morecomment=[s]{\#=}{=\#},
    morestring=[b]",
    morestring=[b]',
}

\lstset{
    language=Julia,
    basicstyle=\ttfamily\small,
    keywordstyle=\color{blue}\bfseries,
    commentstyle=\color{gray}\itshape,
    stringstyle=\color{purple},
    numberstyle=\tiny\color{gray},
    numbers=left,
    stepnumber=1,
    frame=single,
    breaklines=true,
    captionpos=b,
    tabsize=4,
    showstringspaces=false,
}

\lstset{
    literate={ρ}{$\rho$}{1}
             {α}{$\alpha$}{1}
             {β}{$\beta$}{1}
             {σ}{$\sigma$}{1}
             {σ²}{$\sigma^2$}{2}
             {μ}{$\mu$}{1}
             {ϵ}{$\epsilon$}{1}
             {η}{$\eta$}{1}
             {η²}{{$\eta^2$}}{2}
             {r̂}{$\hat r$}{1}
             {Δ}{$\Delta$}{1}
}

\title{STAT 547E - Homework 1}
\author{Zachary Lau}
\date{\today}

\begin{document}
\maketitle
\section{Problem 1}
    % Question 1
    \begin{enumerate}
        \item % 1-a
        To write down the joint kernel we consider the two ways that we can move
        from the state $(x,y)$ to $(x',y')$ these are 
        \begin{itemize}
            \item Draw $x'$ from $K(x,\cdot)$
            \item Draw $y'$ from $\eta(\cdot)$
            \item Reject the swap. This happens with probability
            \begin{align}
                1 - \alpha(x',y') &= 1 - \min\left\{1, \frac{w(y')}{w(x')}\right\} \\ 
                &= 1 - \min\left\{1, \frac{\gamma(y')\eta(x')}{\gamma(x')\eta(y')}\right\}
            \end{align}
        \end{itemize}
        and
        \begin{itemize}
            \item Draw $y'$ from $K(x,\cdot)$
            \item Draw $x'$ from $\eta(\cdot)$
            \item Accept the swap. This happens with probability
            \begin{align}
                \alpha(y', x') &= \min\left\{1, \frac{w(x')}{w(y')} \right\} \\
                &= \min\left\{1, \frac{\gamma(x')\eta(y')}{\gamma(y')\eta(x')} \right\} 
            \end{align}
        \end{itemize}
        Combining the two and noting that the swap is volume preserving
        (so Jacobian determinant is 1), we get the expanded kernel
        \begin{align}
            \bar K \left((x,y),\d(x',y')\right) =
            K(x,\d x')\eta(\d y')
            \left[1 - \min\left\{1, \frac{\gamma(y')\eta(x')}{\gamma(x')\eta(y')}\right\}\right]
            + \\
            K(x,\d y')\eta(\d x')
            \min\left\{1, \frac{\gamma(x')\eta(y')}{\gamma(y')\eta(x')} \right\} 
        \end{align}
        Note that we have slightly abused notation to let $\eta$ refer to both
        a density (in the acceptance weights) and a measure (outside them).
        

        We show invariance directly from the definition. Let 
        $\pi$ refer to the measure with density $\frac{\gamma(x)}{Z}$. We
        see that the joint distribution over $((x, y),(x', y'))$ is given by
        \begin{align}
            \bar \pi ((\d x,\d y))\bar K \left((x,y),(\d x', \d y')\right) &=
            \pi(\d x)\eta(\d y)
            K(x,\d x')\eta(\d y')
            \left[1 - \min\left\{1, \frac{\gamma(y')\eta(x')}{\gamma(x')\eta(y')}\right\}\right] \\
            &+ 
            \pi(\d x)\eta(\d y)K(x,\d y')\eta(\d x')
            \min\left\{1, \frac{\gamma(x')\eta(y')}{\gamma(y')\eta(x')} \right\} 
        \end{align}
        We can integate the summands one at a time (note both are almost-surely
        non-negative). Consider the second summand in the above equation
        \begin{align}
            S_2 = \pi (\d x)\eta(\d y) K(x, \d y') \eta(\d x')  \min\left\{1, \frac{\gamma(x')\eta(y')}{\gamma(y')\eta(x')} \right\} \\
        \end{align}
        We use Fubini (these are all positive meaures)
        and marginalize over $x$ first. Note that $K$ is stationary
        for $\pi$ so $\int \pi(\d x)K(x,\d y') = \pi(\d y')$.
        Furthermore, the only factors depending on $x$ are $\pi(\d x)$ and
        $K(x, \d y')$.
        \begin{align}
            \int_{\mathbb X} S_2 &= 
            \eta(\d y)\eta(\d x') \min\left\{1, \frac{\gamma(x')\eta(y')}{\gamma(y')\eta(x')} \right\} 
            \int \pi(\d x)K(x, \d y') \\ 
            &=  \eta(\d y)\eta(\d x') \min\left\{1, \frac{\gamma(x')\eta(y')}{\gamma(y')\eta(x')} \right\}
                \pi(\d y')
        \end{align} 
        The only factor dependig on $y$ is $\eta (\d y)$ so marginalizing over
        $y$ we get 
        \[
            \int_{\mathbb X \times \mathbb X} S_2 = 
                \eta(\d x') \min\left\{1, \frac{\gamma(x')\eta(y')}{\gamma(y')\eta(x')} \right\}
                \pi(\d y')
        \]
        Noting that $\eta(\d x ') = \eta(x') \d x'$ and
        $\pi (\d y') = \frac{\gamma(y')}{Z}$ we can further simplify
        \begin{align}
            \int_{\mathbb X \times \mathbb X} S_2 =
                \frac 1 Z \min\left\{\gamma(y')\eta(x'), \gamma(x')\eta(y') \right\} \d x' \d y'
        \end{align}
        Similarly for the first term we get
        \begin{align}
            \int_{\mathbb X \times \mathbb X} S_1 &=
                \int_{\mathbb X \times \mathbb X} 
                \pi(\d x)\eta(\d y)K(x,\d x')\eta(\d y')
                \left[1 - \min\left\{1, \frac{\gamma(y')\eta(x')}{\gamma(x')\eta(y')}\right\}\right] \\
            &= \pi(\d x')\eta(\d y')
                        \left[1 - \min\left\{1, \frac{\gamma(y')\eta(x')}{\gamma(x')\eta(y')}\right\}\right] \\
            &= \pi(\d x)\eta(\d y') - \frac 1 Z \min\left\{
                \gamma(x')\eta(y'),\gamma(y')\eta(x') 
            \right\} \d x' \d y'
        \end{align}
        The second term cancels with the above. We get
        \begin{align}
            \int_{\mathbb X \times \mathbb X} S_1 + S_2 &= 
                \pi(\d x')\eta(\d y') = \bar \pi(\d x', \d y')
        \end{align}
        Therefore $\bar K$ is $\bar \pi$-invariant.
    \item Firstly we show the proof of the following proposition
    \begin{proposition}\label{prop:tv}
        Supposing that $\mu$ has density $\mu(x)$ and $\nu$ has density $\nu(x)$,
        then the TV distance $T(\mu, \nu)$ is given by $1 - \int \min(\mu(x),\nu(x)) \d x$
    \end{proposition}
    We use the following Lemma
    \begin{lemma}\label{lemma:tv}
        Given $\mu$ and $\nu$ above then
        $\arg\max_A \nu(A) - \mu(A) = \{ x : \nu(x) > \mu(x) \}$ 
    \end{lemma}
    \begin{proof}[Proof of Lemma \ref{lemma:tv}]
        Let $A = \{ x : \nu(x) > \mu(x) \}$. Given any set $A'$ we see that
            $\nu(A \cap A'^c) \geq \mu(A \cap A'^c)$ and
            $\nu(A^c \cap A') \leq \mu(A^c \cap A')$
        \begin{align}
            \nu(A) - \mu(A) &= 
            \nu(A \cap A') + \nu(A \cap A'^c) 
            - \mu(A \cap A') - \mu(A \cap A'^c) \\
            &\geq \nu(A \cap A') - \mu(A \cap A') \\
            &\geq \nu(A \cap A') - \mu(A \cap A') + 
            \nu(A^c \cap A') - \mu(A^c \cap A') \\
            &= \nu(A') - \mu(A')
        \end{align}
    \end{proof}
    \begin{proof}[Proof of proposition \ref{prop:tv}]
        By Lemma \ref{lemma:tv} there exist maximizers $A$ and $B$ of
        $\nu(\cdot)-\mu(\cdot)$ and $\mu(\cdot) - \nu(\cdot)$ of the form 
        $A = \{ x: \nu(x) > \mu(x) \}$ and $B = \{ x : \nu(x) < \mu(x) \}$.
        At least one of
        these maximizes $|\mu(\cdot) - \nu(\cdot)|$. Suppose that
        $\nu(A) - \mu(A) \geq \mu(B) - \nu(B)$. Then we can calculate the TV
        distance
        \begin{align}
            \TV(\nu,\mu) &= \nu(A) - \mu(A) \\
                &= \int \1\{ \nu(x) > \mu(x) \} \nu(x) \d x
                - \int \1\{ \nu(x) > \mu(x) \} \mu(x) \d x \\
                &= \int \1\{ \nu(x) > \mu(x) \} \nu(x) \d x
                - \left[ 1 - \int \1\{ \nu(x) \leq \mu(x) \} \mu(x) \d x \right] \\
                &= \int \left[\1\{ \nu(x) > \mu(x) \} \nu(x) + \1\{ \nu(x) \leq \mu(x) \} \mu(x)
                \right] \d x - 1 \\
                &= \int \max\{ \nu(x), \mu(x) \} \d x - 1
        \end{align}
        Note that $\max(a,b) = a + b - \min(a,b)$ so we get
        \begin{align}
            \TV(\nu, \mu) &= \int \nu(x) + \mu(x) - \min\{ \mu(x), \nu(x) \} \d x - 1 \\
            &= 1- \int \min(\mu(x), \nu(x)) \d x
        \end{align}
        Integrating over $B$ we get the same result, so our initial assumption does
        not matter, and the lemma is proven.
        \end{proof}
        We return to the immediate problem at hand. Given $(x,y) \sim \bar \pi$,
        the pre-swap proposal is distributed according to $\bar \pi$ as well, thus
        the swap acceptance rate is given by
        \begin{align}
            \int \pi(x)\eta(y)\min\left\{1, \frac{\gamma(y)\eta(x)}{\gamma(x)\eta(y)} \right\}  \d x \d y
            &= \int \pi(x)\eta(y)\min\left\{1, \frac{\pi(y)\eta(x)}{\pi(x)\eta(y)} \right\} \d x \d y  \\
            &= \int \min\left\{ \pi(x)\eta(y), \pi(y)\eta(x)\right\} \d x \d y \\
            &= 1- \TV(\pi \otimes \eta, \eta \otimes \pi) \\
            &= 1 - \|\pi \otimes \eta - \eta \otimes \pi \|_{\TV}
        \end{align}
\end{enumerate}

\FloatBarrier
\newpage
\section{Problem 2}
\begin{enumerate}
    \item % a
    % TODO formalized this
    This proposition is true. We provide a heuristic proof. Suppose that our
    chain starts at $(x_0, y_0)$ and consider any measurable set
    $A \subset \mathbb X \times \mathbb X$ with $\bar \pi (A) > 0$.
    We wish to show that the chain will reach $A$ with positive probability by
    some finite time $t$.
    Let $A_x  = \{ x : \eta\{ y: (x,y) \in A\} > 0 \}$
    (i.e. all $x$ which occur with positive probability in $A$). Then marginalizing over 
    $y$ we see that $\pi(A_X) > 0$. Therefore by the $\pi$-irreducibility of
    $K$ there is a sequence of transitions that will move $x$ to $A_x$ in the
    original non-stabilized chain with positive probability by time $t$. WLOG of generality, we assume that the
    starting point has positive density $\pi(x_0) > 0$, otherwise, the first swap
    will immediately be accepted and the proof proceeds from that point.
    We now argue that such an $x$-trajectory has positive probability even in
    the presence of swaps. Given that
    $\pi(x_0) > 0$, then since $K$ is $\pi$-irreducible each following $x_i$
    will have $\pi(x_i) > 0$. Now for each $i$ either $\eta(x_i) = 0$, in which case the RN
    derivative is infinite and all swaps are rejected, or $\eta(x_i) > 0$ in
    which case there is positive probability that $x_i$ is drawn from the
    reference and the swap leaves the state unchanged. Either way, our $x$
    trajectory is unaffected. All that remains is that at the last step a valid
    $y_t$ is drawn such $(x_t, y_t) \in A$, which occurs with positive probability
    by our definition of $A_x$.

    % Consider any starting
    % point $x$ and set $A \subset \mathbb X$ with $\pi(A) > 0$.
    % By the $\pi$-irreducibility of $K$, for some $t$ we have
    % $\pi K^t(A) > 0$. For such a $t$ consider the space of all such paths i.e.
    % $\nu = \pi(\d x_0) \otimes K(x_0, \d x_1) \otimes \ldots \otimes K(x_{t-1}, \d x_t)$,
    % and let $B$ be the set of all paths which end in $A$. Then 
    % $\nu(B) = \pi K^t (A) > 0$. Since $K$ is $\pi$-irreducible, for each $i$,
    % $w(x_i) > 0$, and thus there is 
    

    % The image of $D$ under $\bar K$ can be found the following procedure
    % \begin{itemize}
    %     \item Take $D'_X$ the image of $D_X$ under $K$
    %     \item $D'_Y$ is the support of $\eta$, since $y$ is drawn from the reference
    %     \item This gives $D' = D'_X \times D'_Y$ (since $y$ is drawn independently of $x$)
    %     \item Compute $E$ the image of $D'$ under the swapping kernel
    %     \[ K((x,y),\d (x,y)) = \alpha(x,y)\delta_{(y,x)}(x',y') + (1-\alpha(x,y)) \delta_{(x,y)}(x', y')\]
    % \end{itemize}
    % We need to show that if $(\pi \otimes \eta)(D) < 1$ then
    % $(\pi \otimes \eta )(E \setminus D) > 0$. 


    % This proposition is true. Supposing that for some set $A \subset X$
    % $P(X_t \in A) = \epsilon > 0$ for some $t < \infty$ in the reduced chain.
    % Then consider the sequence of acceptance probabilities
    %     $\alpha_t =$ 
    \item % b
    This is false. We provide a counter example on a two-state Markov chain
    with states $0$ and $1$. Let the reference be given by
     $\eta(0)=\eta(1)=\frac 1 2$, transition kernel by $K(x,\d x')=\delta_x(\d x')$
     and target by $\pi(0) = \pi(1) = \frac 1 2$.
    $K$ is not $\pi$-irreducible. For example, if a chain starts at $0$ it will be stuck at $0$
    forever and never reach $1$. However, $\bar K$ is $\bar\pi$-irreducible, in-fact
    since $w(0) = w(1) = 1$, swaps are always accepted and so $P(X_T = 1) = P(X_T = 0) = \frac 1 2$
    for all $T > 0$ regardless of starting point. 
    \item % c
    % This is true and the proof is similar to question 1. We can just plug and
    % chug. Assuming $K$ is $\pi$-reversible then $\pi(\d x)K(x, \d x') = \pi(\d x')K(x', \d x)$
    % Plugging this in and cancelling terms we get
    % \begin{align}
    %     &\bar \pi(\d (x,y)) \bar K((x,y), \d(x',y')) \\
    %     &= 
    %             \bar \pi ((\d x,\d y))\bar K \left((x,y),(\d x', \d y')\right) \\
    %             &= 
    %         \pi(\d x)\eta(\d y)
    %         K(x,\d x')\eta(\d y')
    %         \left[1 - \min\left\{1, \frac{\gamma(y')\eta(x')}{\gamma(x')\eta(y')}\right\}\right] 
    %         + 
    %         \pi(\d x)\eta(\d y)K(x,\d y')\eta(\d x')
    %         \min\left\{1, \frac{\gamma(x')\eta(y')}{\gamma(y')\eta(x')} \right\}  \\
    %     &=  
    %         \pi(\d x')\eta(\d y)
    %         K(x',\d x)\eta(\d y')
    %         \left[1 - \min\left\{1, \frac{\gamma(y')\eta(x')}{\gamma(x')\eta(y')}\right\}\right] 
    %         + 
    %         \pi(\d y')\eta(\d y)K(y',\d x')\eta(\d x')
    %         \min\left\{1, \frac{\gamma(x')\eta(y')}{\gamma(y')\eta(x')} \right\}  \\
    %     &= \pi(\d x')\eta(\d y)K(x',\d x)\eta(\d y')
    %     - \frac 1 Z \min\left\{\gamma(x')\eta(y'), \gamma(y')\eta(x') \right\} \d x' \d y'
    %     + \frac 1 Z \min\left\{\gamma(x')\eta(y'), \gamma(y')\eta(x') \right\} \d x' \d y' \\
    %     &= \pi(\d x')\eta(\d y)K(x',\d x)\eta(\d y')
    % \end{align}
    % This expression is symmetric in $y$ and $y'$. We can apply the $\pi$-reversibility of
    % $K$ to swap the $x$ and $x'$ and then simply reverse our previous steps
    % \begin{align}
    %     &= \pi(\d x)\eta(\d y)K(x, \d x')\eta(\d y') \\
    %     &= \bar \pi (\d(x', y')) \bar K((x', y'),\d (x, y))
    % \end{align}
    % Thus $\bar K$ is $\bar \pi$-reversible.
    This statement is false. We provide a counter example on a two-state Markov
    chain. Denoting measures by row vectors let
    \begin{itemize}
        \item $\pi = \left[\frac 1 2, \frac 1 2 \right]$
        \item $\eta = \left[ \frac 2 3, \frac 1 3 \right]$
    \end{itemize}
    Denoting the first state by $0$ and the second state by $1$ we get
    $w(0) = \frac 3 4$ and $w(1) = \frac 3 2$. This gives
    $\alpha(0, 1) = 1$ and $\alpha(1, 0) = \frac 1 2$. The joint distribution
    at stationarity is given by the product of $\pi$ and $\eta$, namely
    $\bar \pi = \left[\frac 1 3, \frac 1 6, \frac 1 3, \frac 1 6\right]$
    , where probabilities correspond to the the states
    $(X=0, Y=0), (X=0, Y=1), (X=1, Y=0)$ and $(X=1,Y=1)$ respectively.
    Let our local kernel be the identity kernel. This kernel is trivially
    reversible with respect to any measure.
    We now consider the detail balance condition for the states $(0, 1)$ and
    $(1,0)$ with this kernel. A transition from state $(0,1)$ to $(1, 0)$ can
    only occur if the proposal $Y$ is $1$ and the swap is accepted. This 
    gives
    \[
        \bar \pi (0, 1)\bar K((0, 1), (1,  0)) = 
            \frac 1 6 \times \frac 1 3 \times  1 = \frac 1 {18}
    \]
    On the other hand from $(1,0)$ to $(0, 1)$ we have
    \[
        \bar \pi(1, 0)\bar K((1, 0), (0, 1)) = 
            \frac 1 3 \times \frac 2 3 \times \frac 1 2 = \frac 1 9
    \]
    These are not equal so detail balance is violated - $\bar K$ is not
    $\bar \pi$-reversible. As a side note, interestingly both the local
    exploration step and swap step are themselves reversible. However, an
    arbitrary combination of reversible kernels is not necessarily reversible,
    as shown in this example.
\end{enumerate}

\FloatBarrier
\newpage
\section{Problem 3}
\begin{enumerate}
    \item % a
    From Problem 2 we know that at stationarity the expected acceptance
    rate is given by
    $\E[\alpha] = 1-\| \pi \otimes \eta - \eta \otimes \pi \|_{\TV}$. 
    Using $\pi = \Norm(\mu, \sigma^2)$ and $\eta = \Norm(0, \sigma^2)$ we
    we get \[ \pi \otimes \eta \sim \Norm([\mu, 0], \sigma^2 I_2)\] and 
    \[ \eta \otimes \pi \sim \Norm([0, \mu], \sigma^2 I_2)\]. Applying the hint
    and with $\Sigma = \sigma^2 I_2$ we get
    \begin{align}
        \|\pi \otimes \eta - \eta \otimes \pi \|_{\TV} &= 
            2\phi\left(\frac 1 2 \| [\mu, -\mu] \|_{\Sigma^{-1}}\right) - 1 \\
            &= 2\phi\left(\frac 1 2 \sqrt{\frac{2\mu^2}{\sigma^2}}\right) - 1 \\
            &= 2\phi\left(\frac{\rho}{\sqrt{2}}\right) - 1
    \end{align}
    We next show that $\frac{1}{T}\sum_{i=1}^T \alpha_t$ converges almost surely
    to $\E_{\bar\pi}[\alpha]$. To do so we simply need to show that the
    chain specified is $\bar\pi$-invariant, aperiodic and $\bar\pi$-irreducible. Invariance
    comes from question 1. Irreducibility follows from the global support
    of $\eta$. Since $\eta$ has infinite support $\eta(y)$ can be arbitrarily
    small ensuring there is always a non-zero chance the swap is rejected, thus
    we have aperiodicity. Finally note that $\pi$ has global support so we can
    apply Theorem 4 of Roberts and Rosenthal \cite{Roberts_2004} for any $x \in \mathbb R$.

    % TODO flesh out irreducibility in more detail
    Given that $\frac{1}{T}\sum_{i=1}^T \alpha_t$ converges almost surely
    to $\E_{\bar\pi}[\alpha]$ we find that 
    $\hat r_T = 1 - \frac{1}{T}\sum_{i=1}^T \alpha_t$ converges almost
    surely to 
    \begin{align}
        1 - \E_{\bar\pi}[\alpha] &= 
            \|\eta - \pi \otimes \pi - \eta\|_{\TV} \\
        &= 2\phi\left(\frac{\rho}{\sqrt{2}}\right) - 1 
    \end{align}

    Listing \ref{lst:wait} implements the wait time as a function of
    $\rho$ in Julia and plots the
    resulting wait times for $\rho \in [0,5]$. The resulting curves are shown in
    Figures \ref{fig:rejection} (with $\rho$ on a linear scale) 
    and \ref{fig:log_rejection} (with $\rho$ on a log-scale). The results show
    that for higher $\rho$ there are more rejections, and we have to go a longer
    time between swaps.

    \begin{lstlisting}[language=Julia, caption={Code for theoretical swap rejeciton
        rates and waiting times}, label={lst:wait}]
function r_fn(ρ)
    2*cdf(Normal(),ρ/sqrt(2)) - 1
end

function wait_fn(ρ)
    r_fn(ρ)/(1-r_fn(ρ))
end

plot_ρs = range(0, 5, length=100)
plot_r = r_fn.(plot_ρs)
plot_wait = wait_fn.(plot_ρs)
plot(
    plot(plot_ρs, plot_r, label="r", color="red", title="Average rejection rate",
        xlabel = "ρ",  ylabel = "Average rejection rate at stationarity"),
    plot(plot_ρs, log10.(plot_wait), label="log10 w", title="Waiting time",
        xlabel = "ρ", ylabel = "Log10 waiting time")
)

plot_ρs = 10 .^ range(-2, 1, length=100)
plot_r = r_fn.(plot_ρs)
plot_wait = wait_fn.(plot_ρs)
plot(
    plot(log10.(plot_ρs), plot_r, label="r", color="red", title="Average rejection rate",
        xlabel = "log10 ρ",  ylabel = "Average rejection rate at stationarity"),
    plot(log10.(plot_ρs), log10.(plot_wait), label="log10 w", title="Waiting time",
        xlabel = "log10 ρ", ylabel = "Log10 waiting time")
)
    \end{lstlisting}
    \FloatBarrier

\begin{figure}
    \centering
    \includegraphics[width=\textwidth]{rejection.png}
    \caption{Theoretical rejection rate and waiting times}
    \label{fig:rejection}
\end{figure}

\begin{figure}
    \centering
    \includegraphics[width=\textwidth]{log_rejection.png}
    \caption{Theoretical rejection rate and waiting times with $\rho$ on a log-scale}
    \label{fig:log_rejection}
\end{figure}

    \item First we \emph{should} show that the $Y_t$ are \iid draws from $\eta$. This is true
    at stationarity. The $Y_t$ are by definition drawn from $\eta$. Furthermore,
    since $\bar \pi = \pi \otimes \eta$, the $X_t$ are independent of $Y_t$. This
    breaks the causal graph between $Y_t$ and $Y_{t+1}$, so they are
    independent.
    
    When we are not at stationarity,
    $Y_t$ are not necessarily drawn from $\eta$. We provide a brief discussion
    of this point. Firstly, we demonstrate this empirically through a simulation of the normal
    example used in this question with $\sigma^2=\mu=1$. A simulation is implemented in 
    Listing \ref{lst:biased} and the output is shown in Figure \ref{fig:biased}.
    We can see that prior to reaching stationarity the $Y$'s are actually skewed
    left as any higher values are more likely to get swapped to $X$'s.

    We can also provide an exact calculation on a two-state Markov chain to
    show that in the short-term the $Y$-values are biased in the stabilized MCMC
    algorithm. We consider the simple case with target $\pi = [\frac 2 3, \frac 1 3]$ on
    state space $\{0,1\}$ and reference $\eta = 
    \left[\frac 1 2, \frac 1 2 \right]$. With states
    ordered as $[(0,0), (0,1), (1,0), (1,1)]$, this gives
    stationary distribution
    $\bar \pi = \left[\frac 1 3, \frac 1 3, \frac 1 6, \frac 1 6 \right]$. We use as our local kernel the
    delta kernel $K(x, \d x') = \delta_x(\d x')$. The
    weights are $w(0) = \frac 4 3 $ and $w(1) = \frac 2 3 $, swap acceptances 
    $\alpha(0,1) = \frac 1 2$ and $\alpha(1,0) = 1$ and hence transition matrix
    \[ K = \begin{bmatrix}
        \frac 1 2 & \frac 1 4 & \frac 1 4 & 0 \\
        \frac 1 2 & \frac 1 4 & \frac 1 4 & 0 \\
        0 & \frac 1 2 & 0 & \frac 1 2 \\ 
        0 & \frac 1 2 & 0 & \frac 1 2
    \end{bmatrix} \]
    
    Initializing our distribution at $\nu_0 = \eta \otimes \eta = 
    \left[\frac 1 4, \frac 1 4, \frac 1 4, \frac 1 4\right]$ we find the distribution of the chain after the first
    few steps is given by 
    \begin{itemize}
        \item $\nu_1 =  [0.2500, 0.3750, 0.1250, 0.2500]$
        \item $\nu_2 = [0.3125, 0.3438, 0.1562, 0.1875]$
        \item $\nu_3 = [0.3281, 0.3359, 0.1641, 0.1719]$
    \end{itemize}
    And specifically looking at the marginal distribution of $Y$ we find
    \begin{itemize}
        \item $Y_1 \sim [0.3750, 0.6250]$
        \item $Y_2 \sim [0.4688, 0.5312]$
        \item $Y_3 \sim [0.4922, 0.5078]$
    \end{itemize}
    Not only is the marginal distribution of $Y_i$ not $\eta$, from the joint
    distributions we can see that $X$ and $Y$ are not independent. E.g. for the 
    first time step $0.25 \times 0.125 \neq 0.375 \times 0.25$, while for
    independent $X$ and $Y$
    $P(0,1)P(1,0)=P(X=0)P(Y=1)P(X=1)P(Y=0)=P(0,0)P(1,1)$.
    Further discussion might include analyis of the combined state markov chain
    with state $S_i = (X_i, X_{i+1}, Y_i, Y_{i+1})$ to show dependence between
    $Y_i$ and $Y_{i+1}$.

    \begin{lstlisting}[language=Julia, caption={Simulation of stabilized MCMC
        before stationarity on the normal example with 
        $\eta \sim \Norm(0,1)$ and $\pi \sim \Norm(1,1)$}, label={lst:biased}]
function logw(x)
    μ = 2
    x*μ-1/2*μ^2 # Target is mean 1, sd 1 reference is mean 0
end

N = 100000
T = 15
spacing = 5
checkpoints = Matrix{Float64}(undef, N, div(T, spacing))
for j in 1:N 
    c = Float64[]
    x = randn()
    y = randn()
    xs = [x]
    ys = [y]
    acc = 0
    for i in 1:T
        y = randn()
        alpha = logw(y) - logw(x)
        if alpha > 0 || rand() < exp(alpha)
            (x, y) = (y, x)
            acc = acc + 1
        end
        push!(xs, x)
        push!(ys, y)
        if i % spacing == 0
            push!(c, y)
        end
    end
    checkpoints[j,:] = c
    # println(mean(xs))
    # println(acc/T)
end

println(mean(checkpoints, dims = 1))
println(std(checkpoints, dims = 1))

data = checkpoints[:,1]
histogram(data, normalize=:pdf, label="Y_$spacing", alpha = 0.5)
density!(data, label="KDE", linewidth=2)
plot!(Normal(0, 1), label="N(0,1)", linewidth=2)
    \end{lstlisting}
\FloatBarrier
\begin{figure}
    \centering
    \includegraphics[width=\textwidth]{biased.png}
    \caption{Distribution of $Y$ after 5 iterations of stabilized MCMC on
    normal target with $\eta \sim \Norm(0,1)$ and 
    $\pi \sim \Norm(1,1)$}
    \label{fig:biased}
\end{figure}
    
    For the sake of
    simplicity we assume that $\hat Z_T$ 
    is defined in terms of the $Y'_t$ which \emph{are} \iid draws from 
    $\eta$. In this case
    \[
        \E w(y) = \int \frac{\pi(y)}{\eta(y)}\eta(y) \d y = 1
    \]
    Next we can compute the variance of the weights.  For
    this specific case, the weights are given by 
    \begin{align}
        \log(w(y)) &= -\frac 1 2 \frac{(y-\mu)^2}{\sigma^2} + \frac 1 2 \frac{y^2}{\sigma^2} \\
        &= \frac{2\mu y - \mu^2}{2\sigma^2} \\
        w(y) &= \exp\left(\frac{2\mu y - \mu^2}{2\sigma^2}\right)
    \end{align}
    Given that $y \sim \Norm(0, \sigma^2)$ we see that 
    $\frac{2\mu y - \mu^2}{2\sigma^2}$ is normal with mean 
    $-\frac{\mu^2}{2\sigma^2} = -\frac 1 2 \rho^2$ and variance
    $\sigma^2 \left( \frac{\mu}{\sigma^2} \right)^2 = \rho^2$. Thus
    $w(y)$ is log-normally distributed with $\log w(y) \sim \Norm\left(-\frac 1 2 \rho^2, \rho^2 \right)$.
    From the log-normal distriubtion properties this gives
        $\E w(y) = 1$ and $\Var w(y)= \exp(\rho^2) - 1$.
    Since the $y$ are \iid the $w(y)$ are also \iid and we find that
    $\E \hat Z_T = 1$ and $\Var \hat Z_T = \frac {\exp(\rho^2) - 1} T$.
    A simulation of the importance sampling weights is done in
    listing \ref{lst:is}. The resulting empirical mean and scaled variance is shown in 
    Table \ref{tab:is}. For small $\rho$ the empirical quantities are close
    to what we expect. For larger $\rho$ due the distance between our
    proposal and target distribution the estimates of both the mean and variance
    are much less reliable.

\begin{lstlisting}[language=Julia, caption={Importance Sampling Estimate of Normalizing Constant},
    label={lst:is}]
function logw(y,ρ)
    ρ*y-1/2*ρ^2 # Cancel the y^2 terms
end

T = 200
y = randn(T)
# We can use common random numbers to reduce variability
for ρ in [0.5, 1, 2, 3]
    println("=== ρ=$ρ ===")
    ls = logw.(y, ρ)
    w = exp.(ls)
    avg = mean(w) # estimate mean
    variance = var(w) # estimate variance
    println("Mean $avg and variance $variance")
    tvar = exp(ρ^2) - 1
    println("Theoretical variance: $tvar")
end
\end{lstlisting}

\begin{table}[h]
\centering
\csvautotabular{is.csv}
\caption{Mean and Variance of Importance Sampling Estimates of Normalizing
Constant. The true value is $1$. Variances are scaled by $T$, the number of
samples.}
\label{tab:is}
\end{table}

    \item % c
    Viewing the renewals as coming from a renewal process,
    the expected waiting time until a renewal at stationarity is given by 
    \begin{align}
        \frac{1}{\E[\alpha]} - 1 &= 
        \frac{r(\rho)}{1 - r(\rho)} \\ &=
            \frac{2\phi\left(\frac{\rho}{\sqrt{2}}\right) - 1}
            {2 -2\phi\left(\frac{\rho}{\sqrt{2}}\right)} \\
    \end{align}

    Near $\rho=0$ the numerator approaches $0$ and denominator approaches $1$.
    We can take a first order Taylor approximation of the normal CDF, recalling
    that it's derivative is the normal pdf $\frac{1}{\sqrt{2\pi}}\exp\left(-\frac 1 2 x^2\right)$.
    Thus near $\rho=0$ the waiting time behaves like
    \[
        w(\rho) \approx \frac{\rho}{2\sqrt{\pi}} + o(\rho^2)
    \]
    I.e. it approaches $0$ linearly.

    Next we anlayze $\Var \hat Z_T$ as $\rho \to 0$. Taking
    a first order taylor expansion of the exponential we see 
    \[ 
        \hat Z_T(\rho) \approx \frac{1}{T}\rho^2 + o(\rho^4)
    \]
    This approaches $0$ quadratically, which is faster than the rate at which
    the wait time approaches zero.

    Next we consider the behaviour of $\Var \hat Z_T$ as $\rho \to \infty$.
    The first term dominates and $\Var \hat Z_T$ is
    $O\left( \exp(\rho^2) \right)$.

    As $\rho \to \infty$, $\phi\left(\frac{\rho}{\sqrt{2}}\right) \to 1$. Thus the
    numerator of $w(p)$ approaches $1$ and the denominator approaches $0$.
    The waiting time thus scales like
    $\frac{1/2}{1-\phi(\rho/\sqrt{2})} = \frac{1}{2S(\rho/\sqrt{2})}$ where $S$ is the normal survival
    function. We can compare this to the rate at which $w(\rho)$ goes to infinity
    by using l'Hopitals. We get
    \begin{align}
        \lim_{\rho\to\infty} \frac{w(\rho)}{\Var \hat Z_T (\rho)} 
        &=  \frac{2}{T}\lim_{\rho\to\infty} \frac{\exp(-\rho^2)}{S\left(\frac{\rho}{\sqrt{2}}\right)}
        \times \frac{1}{1-\exp(-\rho^2)}\\
        &= \lim_{\rho\to\infty} \frac{-2\rho\exp(-\rho^2)}
            {-\frac{1}{\sqrt{2\pi}}\exp\left(-\frac{\rho^2}{4}\right)} \\
        &= 0
    \end{align}
    So the variance of the normalizing constant estimate approaches infinity
    faster than the waiting time.

    Listing \ref{lst:wait_vs_var} plots the waiting time and importance sampling variance.
    The plots are shown in Figure \ref{fig:wait_vs_var} on a log-log scale. We can
    see that the behaviour agrees with what we expect in both extremes; the
    variance of the estimate of the normalizing constant both grows faster as
    $\rho$ grows, and shrinks faster as $\rho$ shrinks (note the difference in
    scales).

\begin{lstlisting}[language=Julia, caption={Plotting of wait time and variance
    of normalizing constant estimator}, label={lst:wait_vs_var}]
function var_t(ρ, T = 1)
    (exp(ρ^2) - 1)/T
end

plot_ρs = 10 .^ range(-2, 1, length=100)
plot_wait = wait_fn.(plot_ρs)
plot_var = var_t.(plot_ρs)
plot(
    plot(log.(plot_ρs), log.(plot_wait), label="log w", color="red"),
    plot(log.(plot_ρs), log.(plot_var), label="log var T"),
)
\end{lstlisting}

\begin{figure}
    \centering
    \includegraphics[width=\textwidth]{wait_vs_var.png}
    \caption{Waiting time vs Variance of Normalizing Constant Estimator as a
    Function of $\rho$}
    \label{fig:wait_vs_var}
\end{figure}
    
\end{enumerate}

\FloatBarrier
\newpage
\section{Problem 4}
Given that the Gaussians have common diagonal variance the target simplifies as
\begin{align}
    \pi(x) 
    &= \frac {1}{40} \sum_{i=1}^{40} \frac{1}{2\pi \times (0.1)^2}
        \exp\left(-\frac{(x_1-\mu_{i1})^2}{0.2}-\frac{(x_2-\mu_{i2})^2}{0.2} \right) \\
    &= \frac{1}{0.8 \pi} \sum_{i=1}^{40} \exp\left(-\frac{(x_1-\mu_{i1})^2}{0.2}-\frac{(x_2-\mu_{i2})^2}{0.2} \right) \\
    \log(\pi(x)) &=
        C + \log\left(\sum_{i=1}^{40} \exp\left(-\frac{(x_1-\mu_{i1})^2}{0.2}-\frac{(x_2-\mu_{i2})^2}{0.2} \right) \right)
\end{align}
Where $C$ is a constant independent of $x$ and the inner component can be
computed using log-sum-exp. Listing \ref{lst:rwm} implements Metropolis-Hastings
with a random walk proposal.

\begin{lstlisting}[language=julia, caption={RWM MH}, label={lst:rwm}]
##GMM-40 specification
const σ2  = 0.1
const Δ      = 3
const μ_grid = [[i*Δ, j*Δ] for i in 0:4 for j in 0:7]   # 40 means
const K      = length(μ_grid)

function log_gamma(x::Vector)
    logsumexp([-dot(x-μ, x-μ)/(2σ2) for μ in μ_grid])
end

function grad_log_gamma(x::Vector)
    factors = [-dot(x-μ, x-μ)/(2σ2) for μ in μ_grid]
    s = logsumexp(factors)
    sum([(μ-x)/σ2 for μ in μ_grid] .* exp.(factors .- s))
end

# Contour grid for visualisation
xg = range(-2, 14, length=150)
yg = range(-2, 23, length=150)
Z_contour = [exp(log_gamma([x, y])) for y in yg, x in xg];

function rwm_kernel(x, σ_prop; rng)
    # Instrumented rwm kernel (returns acceptances as well)
    # Make proposal
    xp = x + σ_prop*randn(rng, length(x)) # isotropic steps
    # Symmetric proposal so the qs cancel
    alpha = log_gamma(xp)-log_gamma(x)
    acc = false
    if alpha > 0 || rand(rng) < exp(alpha)
        x = xp # else x = x
        acc = true
    end
    (x, acc)
end

function mh_chain(x0, kernel, T; rng)
    # Instrumented mh chain
    acceptances = 0
    x = x0
    chain = Matrix{Float64}(undef, 2, T)
    chain[:,1] = x
    for i in 2:T
        (x, acc) = kernel(x; rng = rng)
        chain[:,i] = x
        acceptances += Int(acc)
    end
    (chain, acceptances/T)
end

# rng is a little overloaded here
rwm_chain(x0, σ_prop, T; rng) = mh_chain(x0,
    (x; rng) -> rwm_kernel(x, σ_prop; rng = rng),
    T ; rng = rng
)
\end{lstlisting}

The gradient is given by 
\[
    \grad \log \pi(x) &= 
    \frac{1}{\sum_{i=1}^{40}\exp\left(-\frac{(x_1-\mu_{i1})^2}{0.2}-\frac{(x_2-\mu_{i2})^2}{0.2}\right)}
    \sum_{i=1}^{40} \begin{bmatrix}
        -\frac{(x_1-\mu_{i1})}{0.1} \\
        -\frac{(x_2-\mu_{i1})}{0.1} 
    \end{bmatrix}
    \exp\left(-\frac{(x_1-\mu_{i1})^2}{0.2}-\frac{(x_2-\mu_{i2})^2}{0.2}\right)
\]
Using our MH infrastructure, and this gradietn we can implement MALA. This is
done in Listing \ref{lst:mala}.

\begin{lstlisting}[language=Julia, caption={MALA}, label={lst:mala}]
function mala_kernel(x, ϵ; rng)
    μ_f = x + ϵ^2/2*grad_log_gamma(x) # Forward proposal
    xp = μ_f + ϵ*randn(rng, length(x)) # Add noise
    μ_b = xp + ϵ^2/2*grad_log_gamma(xp)

    # Log density of proposal is just -(x-μ)^2/(2σ^2)
    df = xp-μ_f
    db = x-μ_b
    alpha = log_gamma(xp) - log_gamma(x) - dot(df,df)/(2ϵ^2) +  dot(db,db)/(2ϵ^2)
    acc = false
    if alpha > 0 || rand(rng) < exp(alpha)
        x = xp
        acc = true
    end
    (x, acc)
end

mala_chain(x0, ϵ, T; rng) = mh_chain(
    x0, (x; rng) -> mala_kernel(x, ϵ; rng = rng), T; rng=Random.default_rng(42)
)
\end{lstlisting}

Both are run for $10000$ iterations with the GMM-40 target in Listing
\ref{lst:run_rwm_mala}.

\begin{lstlisting}[language=Julia, caption={Running RWM and MALA}, label={lst:run_rwm_mala}]
T  = 10_000
x0 = μ_grid[1] .+ 0.1*randn(2)

seed=123

chains_rwm = Matrix{Float64}[]
acc_rwm = Float64[]
σs = [0.3, 1.0, 3.0]
for σ_prop in σs
    println("Running RWM with σ_prop = $σ_prop")
    chain, acceptance = rwm_chain(x0, σ_prop, T; rng=Random.default_rng(42))
    push!(chains_rwm, chain)
    push!(acc_rwm, acceptance)
end

chains_mala = Matrix{Float64}[]
acc_mala = Float64[]
ϵs = [0.05, 0.1, 0.5]
for ϵ in ϵs
    println("Running MALA with ϵ = $ϵ")
    chain, acceptance = mala_chain(x0, ϵ, T; rng=Random.default_rng(42))
    push!(chains_mala, chain)
    push!(acc_mala, acceptance)
end
\end{lstlisting}

The plotting code in Listing \ref{lst:plot_mc} visualizes the results. These results
are shown in Figures \ref{fig:rwm} and \ref{fig:mala} with ESS and acceptance
rates shown in Tables \ref{tab:rwm_gmm40} and \ref{tab:mala_gmm40}. With small
proposal scales, both samplers get stuck in local modes. They have high
acceptance rates, ESS is high, and traceplots look like they have converged.
However, based on the contour plots we can see that they miss many important
modes!

With larger proposal scales, acceptance rates are worse and ess is lower,
but this is not necessarily bad. Larger jumps may be accepted less often, but
this allows us to explore modes further away. In our case, the ESS may be lower, but this is
simply because the ESS is more realistic, and visual analysis of the traceplots
shows that the larger proposal scales explore the space better. Generally, there is a
balance to be had between efficient local exploration (small step sizes) and
the ability to find distance mode (with large step sizes). In this case, the
difficulty in the problem is in the disparate modes, so larger step-sizes win out.

Given that the parameters used are
different for MALA and RWM, we cannot directly compare the two algorithms to one
another, however MALA appears to get even more stuck in local modes than RWM.
Neither sampler is able to convincingly explore all modes in this problem. Even
with larger step sizes it still takes a while to break out and find another mode.

\begin{lstlisting}[language=Julia, caption={Plotting code for MCMC}, label={lst:plot_mc}]
function myplot(alg, chains, parname, values)
    plots = Plots.Plot[]
    for (c, v) in zip(chains, values)
        # Row of scatter plots
        p = contour(xg, yg, Z_contour, levels=3, fill=false, title="$alg with $parname=$v")
        scatter!(p, c[1,:], c[2,:])
        push!(plots, p)
    end
    for (c, v) in zip(chains, values)
        # Row of trace plots
        t = plot(c', alpha=0.5)
        push!(plots, t)
    end
    return plots
end

plot(myplot("RWM", chains_rwm, "σ_prop", σs)..., layout=(2,3), size=(1000,1000))
plot(myplot("MALA", chains_mala, "ϵ", ϵs)..., layout=(2,3), size=(1000,1000))
\end{lstlisting}
\FloatBarrier

\begin{figure}
    \centering
    \includegraphics[width=\textwidth]{rwm.png}
    \caption{Samples from RWM sampler without stabilization}
    \label{fig:rwm}
\end{figure}

\begin{figure}
    \centering
    \includegraphics[width=\textwidth]{mala.png}
    \caption{Samples from MALA sampler without stabilization}
    \label{fig:mala}
\end{figure}

\begin{table}
    \centering
    \csvautotabular{rwm.csv}
    \caption{ESS and Acceptance Ratios for RWM on GMM-40}
    \label{tab:rwm_gmm40}
\end{table}

\begin{table}
    \centering
    \csvautotabular{mala.csv}
    \caption{ESS and Acceptance Ratios for MALA on GMM-40}
    \label{tab:mala_gmm40}
\end{table}

\FloatBarrier
\newpage
\section{Problem 5}
Listing \ref{lst:stable} implements stabilized MCMC.

\begin{lstlisting}[language=Julia, caption={Stabilized MCMC}, label={lst:stable}]
function logw(x, σ_η²)
    # Weighting function for GMM-40 target with normal reference
    return log_gamma(x) + 1/2*dot(x,x)/σ_η²
end

function stabilised_mcmc(x0, σ_η, T; kernel, rng)
    # Note - adding an x0 here as well
    Xs = Matrix{Float64}(undef, 2, T);
    x = x0
    Xs[:,1] = x0
    acceptances = 0
    for i in 2:T
        (x, _) = kernel(x, rng = rng) # Throw away the MH acceptance info
        y = σ_η * randn(length(x0)) # Draw from isotropic reference
        alpha = logw(y, σ_η^2) - logw(x, σ_η^2)
        if alpha > 0 || rand(rng) < exp(alpha)
            x = y # the y = x doesn't matter in practice
            acceptances += 1
        end
        Xs[:,i] = x
    end
    return Xs, acceptances/T
end
\end{lstlisting}

We run this for both RWM and MALA in listing \ref{lst:run_stable} with
$\sigma_\eta = 10$. Resulting
scatter plots are shown in Figures \ref{fig:stabilized_rwm} and 
\ref{fig:stabilized_mala}. The ESS and swap rejection rates are give in 
Table \ref{tab:original}.

\begin{lstlisting}[language=Julia, caption={Running Stabilized MCMC}, label={lst:run_stable}]
σ_η = 10.0
T   = 10_000

# TODO: run classical and stabilised versions; print ESS and rejection rate table
stable_rwm, acc_rwm = stabilised_mcmc(
    x0, σ_η, T;
    kernel = (x; rng) -> rwm_kernel(x, 1, rng = rng),
    rng = Random.default_rng(42)
)

stable_mala, acc_mala = stabilised_mcmc(
    x0, σ_η, T;
    kernel = (x; rng) -> mala_kernel(x, 0.1, rng = rng),
    rng = Random.default_rng(42)
)

original_rwm = chains_rwm[1]
original_mala = chains_mala[1]

function comp_plot(chains, titles)
    plots = Plots.Plot[]
    for (c, title) in zip(chains, titles)
        p = contour(xg, yg, Z_contour, levels=3, fill=false, title=title)
        scatter!(p, c[1,:], c[2,:])
        push!(plots, p)
    end
    for c in chains
        t = plot(c', alpha=0.5)
        push!(plots, t)
    end
    return plots
end

ess_rwm = ess_1d(stable_rwm[1,:])
ess_mala = ess_1d(stable_mala[1,:])
r_rwm = 1-acc_rwm
r_mala = 1-acc_mala
println("RWM: ESS = $ess_rwm, Rejection rate = $r_rwm")
println("MALA: ESS = $ess_mala, Rejection rate = $r_mala")
\end{lstlisting}

\begin{figure}
    \centering
    \includegraphics[width=\textwidth]{stabilized_rwm.png}
    \caption{Samples from stabilized RWM}
    \label{fig:stabilized_rwm}
\end{figure}

\begin{figure}
    \centering
    \includegraphics[width=\textwidth]{stabilized_mala.png}
    \caption{Samples from stabilized MALA}
    \label{fig:stabilized_mala}
\end{figure}

\begin{table}
\centering
\csvautotabular{original.csv}
\caption{ESS of $x_1$ and Swap Rejection Rate for stabilized RWM and MALA
with $\sigma_\eta = 10$}
\label{tab:original}
\end{table}

The stabilized MCMC performs much better than the original. The swap rejection
rate is somewhat high, but it is still able to discover most modes, whereas 
the corresponding unstabilized version is stuck in the initial mode.
% \item 
We repeat for $\sigma_\eta \in \{2^{-5}, 2^{-4}, \ldots, 2^5\}$ in Listing
\ref{lst:compare_stable}. The resulting ESS and empirical rejection rates are
shown in Figures \ref{fig:ess_rwm} and \ref{fig:ess_mala}. The behaviour is
slightly different. For RWM there appears to be a sweet splot near
$\sigma_\eta = 2^4$ where ESS peaks, and the rejection rate dips. However,
MALA actually sees ESS peak near $\sigma_\eta = 2^{-2}$. This is likely to
correspond to a local mode. It is interesting that the $\hat r$ is different
between the two samplers, since they are targetting the same distribution.
This may be due to the fact that we have not yet reached stationarity, and thus
the modes that are actually explord have an outsized effect on $\hat r$.
% A more spread out reference distribution has a higher rejection rate and lower
% ESS, but as we have seen, these are not always useful metrics when talking about
% multi-modal distributions. In particlar, ESS may be high if our chain is stuck
% in a local mode and never explores. This can happen if our referece is itself mostly
% contained in one mode. On the other hand, a broad reference helps us to explore other modes
% better, however the apparent ESS will be lower due to autocorrelation within the
% local mode. Furthermore, the tendency for a broad reference to also include
% areas of low $\pi$-probability increases the swap rejection rate. As we have
% seen from the scatter plots, overall the stabilised version does a much better
% job of discovering all of our modes than the original RWM and MALA. ESS may be
% lower, but in this context it is meaningless since we have yet to reach
% stationarity.

\begin{lstlisting}[language=Julia, caption={Comparing for different values of $\sigma_\eta$}, 
    label={lst:compare_stable}]
σ_η_vals = [2.0^k for k in -5:5]
stable_chains_rwm = Matrix{Float64}[]
stable_chains_mala = Matrix{Float64}[]
as_rwm = Float64[]
as_mala = Float64[]
for σ_η in σ_η_vals
    c_rwm, a_rwm = stabilised_mcmc(
        x0, σ_η, T;
        kernel = (x; rng) -> rwm_kernel(x, 1, rng = rng),
        rng = Random.default_rng(42)
    )
    c_mala, a_mala = stabilised_mcmc(
        x0, σ_η, T;
        kernel = (x; rng) -> mala_kernel(x, 0.1, rng = rng),
        rng = Random.default_rng(42)
    )
    push!(stable_chains_rwm, c_rwm)
    push!(stable_chains_mala,c_mala)
    push!(as_rwm, a_rwm)
    push!(as_mala, a_mala)
end

esses_rwm = [ess_1d(X[1,:]) for X in stable_chains_rwm]
esses_mala = [ess_1d(X[1,:]) for X in stable_chains_mala]

plot(
    scatter(log2.(σ_η_vals), esses_rwm, title="ESS RWM", xlabel="log2(σ_η)", ylabel="ESS"), 
    scatter(log2.(σ_η_vals), 1 .- as_rwm, title="RWM r̂", xlabel="log2(σ_η)", ylabel="r̂")
)
\end{lstlisting}

\begin{figure}
    \centering
    \includegraphics[width=\textwidth]{ess_rwm.png}
    \caption{ESS and empirical rejection rate as a function of $\sigma_\eta$ for
    stabilized RWM}
    \label{fig:ess_rwm}
\end{figure}

\begin{figure}
    \centering
    \includegraphics[width=\textwidth]{ess_mala.png}
    \caption{ESS and empirical rejection rate as a function of $\sigma_\eta$ for
    stabilized MALA}
    \label{fig:ess_mala}
\end{figure}

\bibliographystyle{plain}
\bibliography{refs}

\end{document}